\section{模型的检验}

本节从误差来源、收敛性、参数与口径灵敏度、以及不变量四个方面检验前节的数值结果。
检验所依据的全部数值均由各问的运行记录直接给出，未作任何人工调整。

\subsection{误差来源分析}

数值解的误差可以分解为四项：空间离散误差、时间离散误差、解析近似误差（物性经验式本身的不确定性）、
以及与独立实现比对时暴露的实现差异。前三项通过加密实验与解析对照评估，第四项通过第二套实现交叉验证评估。
需要强调的是，本题的\textbf{主要不确定性并不来自数值误差}，而来自题目给定的物性与环境口径——
这一点在灵敏度分析中被定量地证实。

\subsection{收敛性检验}

对每一问分别做空间与时间的独立加密：空间把网格数由 $400$ 加密到 $800$，时间把内部步长减半或把求解器容差收紧一个量级。
加密后题目要求的位置上的四位小数不再变化，说明所报数值已达到网格与步长无关的水平。
此外，在常物性、固定半径的极限下，本文的实现与圆柱扩散问题的 Bessel 级数解析解对照，
四问的最大绝对误差分别达到 
\qoneanalyticalmaxabserr、\qtwoanalyticalmaxabserr（问题三与其同模型）、\qtwoanalyticalmaxabserr（问题三同模型） 与 \qfourreferencemaxabserr。

需要说明的是，问题一使用的是自适应时间积分，不存在固定步长，因此其时间精度的证据是\textbf{求解器容差加密}
（收紧一个量级后报告点上的表值变化为 \qonetimeaccuracytolerancediff）；
第 (2)(3)(4) 问使用固定步长，证据是步长减半后\textbf{烘干时长}的变化，分别为 
\qtwotimerefinementdiff、\qthreetimerefinementdiff 与 \qfourtimerefinementhours。
问题四另需检验半径插值口径的影响：把附件 2 的单调插值换成线性插值，烘干时长仅变化 \qfourradiusinterphours。

\subsection{参数与口径灵敏度分析}

灵敏度分析针对两类对象：一类是题目给定但数值本身带不确定性的参数（对流传热与传质系数、扩散系数），
另一类是建模过程中必须做出选择的\textbf{口径}（环境外推方式、质量方程的读法、半径插值方式）。
在恒温干燥的基准工况下，各单项扰动对烘干时长的影响汇总于表~\ref{tab:sens}。

\begin{table}[H]
  \centering
  \caption{单项扰动对烘干时长的影响（基准工况的烘干时长为 \qtwodryingtime）}
  \label{tab:sens}
  \small
  \begin{tabular}{lccc}
    \toprule
    扰动项 & 问题二与问题三 & 问题四 & 结论 \\
    \midrule
    对流传质系数放大 $5\%$ & \qtwosenshmplusfivepct & \qfoursenshmplusfivepcthours & 影响很小 \\
    扩散系数放大 $5\%$ & \qtwosensDplusfivepct & \qfoursensDplusfivepcthours & 中等影响 \\
    环境温度升高 $2^\circ$C & \qtwosensTinfplustwoC & \qfoursensTinfplustwoChours & 较大影响 \\
    \bottomrule
  \end{tabular}
\end{table}

表~\ref{tab:sens} 给出的是一个反直觉但对结论至关重要的结果：\textbf{烘干时长对表面传质系数并不敏感}
（在 $5\%$ 的扰动下只变化半个百分点），原因是本题的干燥由\textbf{内部水分扩散控制}，
传质 Biot 数约为 $20$，表面传质阻力不是瓶颈；
因此本文的所有时长结论都写成条件式：\textbf{在题目给定的经验式与质量方程读法之下}。
问题一的规律不同：预热段的表面值对传质系数敏感，放大 $5\%$ 即使表面含水率变化 
\qonehmplusfivepctsurfaceshift，故该问的表面值只在题目给定值下成立。

\subsection{不变量与守恒性检验}

离散格式的守恒性不依赖网格加密，可以\textbf{逐时步核验}。本文把质量守恒写成一个离散恒等式，
四问的残差分别为 \qonemassbalancerel、\qtwomassbalancerel、\qthreemassbalancerel 与 \qfourmassbalancerel，
均达到机器精度量级。

对问题四的移动边界，本文另设一条\textbf{解析不变量检验}：关闭表面传质、只让计算域收缩，
干基含水率必须严格保持不变。若表观对流项的符号或系数写错，收缩就会凭空产生或消灭水分，
该检验因此可以直接\textbf{证伪}项本身的正确性。实测的最大偏离仅为 \qfourmovingboundarydev。

\subsection{与独立实现的交叉验证}

每一问都用一套空间离散与时间推进都不同的实现独立复算：基线采用格心有限体积，表面值由对流边界关系回算，
问题二与问题三的基线还不使用启动步。两套实现在题目要求的报告点上的最大差异为 
\qoneindependentimplementationsdiff（问题一）、\qtwodryingtimespread（问题二与问题三的烘干时长）、
以及 \qfourbaselinetablesixdiff 与 \qfourbaselinedryingtimeseconds（问题四的表值与时长）。

问题二与问题三的时长差异值得单独说明：两套实现在报告点上一致到四位小数的舍入边界，
却仍带来数十秒的烘干时刻差。这不是实现缺陷，而是\textbf{判据位于含水率曲线平坦尾部}的必然结果——
曲线在判据附近近乎水平，极小的数值差异即可换来可观的时间差。因此本文在给出烘干时长的同时，
始终给出该量级的不确定度。

\subsection{跨问一致性}

问题二与问题三的模型、参数与判据完全相同，因此两问应当给出同一个烘干时长。
实测两问分别为 \qtwodryingtime 与 \qthreedryingtime，完全一致（问题三的对照量为 \qthreeconsistencywithQtwo），
这构成一次\textbf{跨问自检}。问题三判据穿越前的最后一个采样点中心含水率为 \qthreeCcenterfivefourh，
穿越后中心与表面分别为 \qthreeCcenteratend 与 \qthreeCsurfaceatend，说明穿越是一次性的、不存在反复越界。

\subsection{检验结论}

四类证据一致表明：本文报告的数值已达到所要求的精度，且其不确定性的主要来源已被识别并量化。
表~\ref{tab:verify} 汇总了四问的全部检验量，图~\ref{fig:verify} 给出四类证据的图形汇总。

\begin{table}[H]
  \centering
  \caption{四问的检验量汇总}\label{tab:verify}
  \footnotesize
  \begin{tabularx}{\textwidth}{l *{5}{>{\centering\arraybackslash}X}}
    \toprule
    问题 & 解析对照/(kg/kg) & 网格加密/(kg/kg) & 时间加密 & 独立实现差 & 离散守恒 \\
    \midrule
    第 (1) 问 & \qoneanalyticalmaxabserr & \qonegridrefinementdiff & \qonetimeaccuracytolerancediff & \qoneindependentimplementationsdiff & \qonemassbalancerel \\
    第 (2) 问 & \qtwoanalyticalmaxabserr & \qtwogridrefinementdiff & \qtwotimerefinementdiff & \qtwodryingtimespread & \qtwomassbalancerel \\
    第 (3) 问 & \qthreedryingtime & \qthreegridrefinementdiff & \qthreetimerefinementdiff & \qthreeconsistencywithQtwo & \qthreemassbalancerel \\
    第 (4) 问 & \qfourreferencemaxabserr & \qfourgridrefinementdiff & \qfourtimerefinementhours & \qfourbaselinedryingtimeseconds & \qfourmassbalancerel \\
    \bottomrule
  \end{tabularx}
  \par\smallskip\footnotesize 第 (3) 问的解析对照取与其同模型的第 (2) 问结果；
  第 (4) 问的时间加密一列为内部步长减半引起的烘干时长变化。
\end{table}


\begin{figure}[H]
  \centering
  \includegraphics[width=0.92\textwidth]{fig05_numerical_verification.pdf}
  \caption{数值可信度的四类证据。(a) 网格加密残差与解析解对照残差；
  (b) 问题四的冻结半径算例与 Bessel 级数解析解的对照；
  (c) 主方法与独立实现的逐点最大差，虚线为一个末位单位对应的容差；
  (d) 问题四关闭表面传质后干基含水率的偏移量与其阈值。}
  \label{fig:verify}
\end{figure}